%% ============================================================
%% 分片 36：原文第 12 节「Branes From Gauge Theory」（原文 glshort.tex 起于第 11186 行）
%% 由 tools/harvmac2tex.py 从 arXiv 源码（hep-th/0604151v3）自动生成。
%% 只译英文散文；公式、编号、\label、\cite 一律不改。
%% ============================================================

\subsection{Branes Of Type $(B,A,A)$}

We can get branes of type $(B,A,A)$ by applying $S$-duality to
branes of type $(B,B,B)$.  Though it may be in general unclear how
to transform a particular brane under $S$-duality, there is no
problem in transforming the unbroken supersymmetries under
$S$-duality.  In view of eqn. \eqref{eq:tokio}, under the transformation
$S$, for ${\rm Re}\,\tau=0$, the generators $\epsilon_{\ell,r}$ of
supersymmetry transform as % 式 (12.20)
\setcounter{equation}{19}
\begin{equation}\label{eq:bilj}
\epsilon_\ell\to
\epsilon_\ell\cdot\exp\left(-{i\pi\over 4}\right),\quad
\epsilon_r\to \epsilon_r\cdot\exp\left({i\pi\over 4}\right),
\end{equation}
 or
equivalently % 式 (12.21)
\setcounter{equation}{20}
\begin{equation}\label{eq:yogo}
\epsilon\to\epsilon'={1+\HG\over \sqrt
2}\epsilon
\end{equation}
 where $\HG=\Gamma_0\Gamma_1\Gamma_2\Gamma_3$ measures
the four-dimensional chirality.  If $T\epsilon=\epsilon$, then
$\epsilon'={1+\HG\over \sqrt 2}\epsilon={1+\HG\over \sqrt
2}T\epsilon={1+\HG\over \sqrt 2}T{1-\HG\over \sqrt 2}\epsilon'$,
so $T'\epsilon'=\epsilon'$ where % 式 (12.22)
\setcounter{equation}{21}
\begin{equation}\label{eq:uvuf}
T'={1+\HG\over \sqrt
2}T{1-\HG\over \sqrt 2}.
\end{equation}
 If $T=\Gamma_0\Gamma_4$, then
% 式 (12.23)
\setcounter{equation}{22}
\begin{equation}\label{eq:polyo}
T'=\Gamma_1\Gamma_2\Gamma_3\Gamma_4.
\end{equation}

We impose the boundary condition $T'\lambda|=-\lambda|$ on the
fermions. Then supersymmetry requires that
$\{T',\Gamma^{IJ}F_{IJ}\}\epsilon=0$, or
% 式 (12.24)
\setcounter{equation}{23}
\begin{equation}\label{eq:bngo}
\left(\sum_{I,J\in
\{1,2,3,4\}}\Gamma^{IJ}F_{IJ}+\sum_{I,J\notin
\{1,2,3,4\}}\G^{IJ}F_{IJ}\right)\epsilon=0.
\end{equation}

This condition implies that $F_{0I}|=0$ for $I>4$, so we impose
Dirichlet boundary conditions on the $A_I$ of $I>4$.  We simply
set $A_I\vert=0$ for $I=5,8,9$, and we look for a boundary
condition that gives a nonzero value for the ``Higgs fields''
$\phi_2=A_6,$ $\phi_3=A_7$.  Assuming that \eqref{eq:bngo}\ is supposed to
be true for all $\epsilon$ obeying
$\Gamma_{1234}\epsilon=\epsilon$ (without the further condition
$\G_{2367}\epsilon=\epsilon$), it implies that
$[\phi_2,\phi_3]=F_{67}=0$.  Given this, the gauge-invariant
content of $\phi_2$ and $\phi_3$ is in the characteristic
polynomial of $\phi_z={1\over 2}(\phi_2-i\phi_3)$.   We pick
Dirichlet boundary conditions in which this characteristic
polynomial is specified.

On the other hand, $A_{2}$, $A_{3}$, and $A_{4}$ obey modified
Neumann boundary conditions % 式 (12.25)
\setcounter{equation}{24}
\begin{equation}\label{eq:yimgo}
\begin{aligned} F_{12}\vert -
F_{34}\vert
 &=0\\
 F_{13}\vert - F_{42}\vert
 &=0\\
 F_{14}\vert - F_{23}\vert
 &=0,\\\end{aligned}
\end{equation}

 which follow from imposing \eqref{eq:bngo}\ with
 $\Gamma_{1234}\epsilon=\epsilon$.

We also want to require the supersymmetries which satisfy
$\Gamma_{2367}\eps=\eps$ to be unbroken. This puts constraints on
the boundary values of the fields which satisfy Dirichlet
conditions. Requiring the vanishing of the supersymmetry
variations of the fermions on the boundary, we find
% 式 (12.26)
\setcounter{equation}{25}
\begin{equation}\label{eq:dirbaa}
\begin{aligned}F_{26}|-F_{73}|&=0,\\ F_{27}|-F_{36}|&=0.\end{aligned}
\end{equation}

These equations are equivalent to $D_\bz\phi_z=0$, which is the
``complex'' Hitchin equation.

These boundary conditions mean that in the Higgs bundle $E$
determined by a pair $(A_{\bar z},\phi_z)$, the field $A_{\bar z}$
which determines the holomorphic structure of $E$ is allowed to
fluctuate, while keeping fixed the characteristic polynomial of
$\phi_z$.  Presumably, the corresponding brane is supported on a
fiber of the Hitchin map $\pi:\MH\to {\EUBB}$.  This gives us an
archetypical $(B,A,A)$-brane.

The condition $[\phi_2,\phi_3]=0$ that arose in this classical
analysis is not one of the Hitchin equations defining $\MH$.
Instead, we should have $[\phi_2,\phi_3]-F_{23}=0$. Presumably,
the extra term arises as a sort of quantum correction in the
renormalization group flow to the infrared.

In the limit that the gauge theory reduces to the two-dimensional
sigma-model, $S$-duality maps a $(B,A,A)$-brane wrapping the fiber
of the Hitchin fibration to a $(B,B,B)$ zerobrane. In the abelian
case, one can verify directly that the boundary conditions that we
have described for these branes are $S$-dual.

We can construct another $(B,A,A)$-brane that seemingly is
likewise supported on a fiber of the Hitchin fibration by
replacing $T'$ with $T''=T'\G_{2367}=-\G_{1467}$.  Imposing
$T''\epsilon=\epsilon$ and $T''\lambda\vert=\lambda\vert$, we get
% 式 (12.27)
\setcounter{equation}{26}
\begin{equation}\label{eq:loppo}
[\G^{IJ}F_{IJ},T'']\epsilon=0,
\end{equation}
 which implies again
Dirichlet boundary conditions for $A_6=\phi_2$ and $A_7=\phi_3$
and Neumann boundary conditions for $A_2$ and $A_3$.

Of course, for a given fiber ${\CMF}$ of the Hitchin fibration
$\pi:\MH\to{{\cmmib B}}$, there should be many branes, associated
with the possible flat line bundles on ${\CMF}$.  We have given
two constructions, but because of time-reversal symmetry, each
leads to a trivial Chan-Paton line bundle (or at most one that is
of order two).  It hopefully is possible to modify the gauge
theory construction to make this bundle vary.

\subsection{Branes Of Type $(A,B,A)$}\label{sec:xgoto}

Our next goal is to construct from the gauge theory some branes of
type $(A,B,A)$. The most interesting of these is the canonical
coisotropic brane. As a prelude, we introduce the abbreviation
$\Gamma_{\pm I\pm J}=(\pm \G_I\pm \G_J)/\sqrt 2$. Similarly, we
take $F_{\pm I\pm J,K}=(\pm F_{IK}\pm F_{JK})/\sqrt 2$, etc.

Let us define % 式 (12.28)
\setcounter{equation}{27}
\begin{equation}\label{eq:jokko}
T'=-\Gamma_{15}\Gamma_{2-7}\Gamma_{3+6}.
\end{equation}

As we will see shortly, this choice leads to branes of type
$(A,B,A)$.  One point of this definition is that $T'$ is invariant
under simultaneous rotation of the $x^2-x^3$ and $x^6-x^7$ planes.
Thus in the twisted theory, $T'$ does not depend on the choice of
the coordinates $x^2,x^3$, provided they are orthogonal. For the
same reason, if we consider the twisted theory on $M=\RR_+^2\times
C$, $T'$ depends only on the metric on $C$, not on the specific
coordinate choice. The branes we will construct using $T'$ depend
only on the conformal class of the metric, i.e. the complex
structure on $C$.

We can again construct two kinds of branes, with boundary
conditions $T'\lambda|=\lambda|$ or $T'\lambda|=-\lambda|$.  In
the former case, supersymmetry gives in the usual way
$[\G^{IJ}F_{IJ},T']\epsilon=0$.  This tells us that $F_{IJ}=0$ if
one of $I,J$ is in the space generated by the $1,5,2-7,$ and $3+6$
directions and the other is orthogonal to it. In particular,
$F_{0I}|=0$ for $I=5,2-7,3+6$, and we impose Dirichlet boundary
conditions for such values of $I$.  But $F_{1I}=0$ for
$I=2+7,3-6,4,8,9$, and we impose Neumann boundary conditions on
$A_I$ for such $I$.

In such a brane, $A_2-\phi_3$ and $A_3+\phi_2$ have specified
boundary values, while $A_2+\phi_3$ and $A_3-\phi_2$ are free to
fluctuate.  To understand this result, it helps to recall that in
the complex structure $J$, the holomorphic coordinates are
$2(A_z+i\phi_z)=(A_2-i A_3)+i(\phi_2-i\phi_3)=(A_2+\phi_3)-i(
A_3-\phi_2)$ and $2(A_{\bar z}+i\phi_{\bar z})=(A_2+i
A_3)+i(\phi_2+i\phi_3)=(A_2-\phi_3)+i( A_3+\phi_2)$.  So our
result means that for this type of brane, ${\cal A}_{\bar
z}=A_{\bar z}+i\phi_{\bar z}$ is fixed and ${\cal A}_{ z}=A_{
z}+i\phi_{ z}$ is unconstrained.  This defines a submanifold $W'$
of $\MH$ which is clearly a complex submanifold in the complex
structure $J$.  In addition, $W'$ is Lagrangian from the point of
view of the holomorphic symplectic form
$\Omega_J=(-i/4\pi)\int_C\Tr\, \delta {\cal A}\wedge \delta {\cal
A}$.   Indeed, $W'$ is middle-dimensional, and in addition
$\Omega_J$ vanishes on $W'$, since $W'$ is characterized by
$\delta {\cal A}_{\bar z}=0$. So the brane supported on $W'$, with
a trivial Chan-Paton bundle, is an $(A,B,A)$-brane, as expected.
$W'$ can be characterized as the submanifold of $\MH$ consisting
of flat $G_{{\Bbb{C}}}$ bundles with a specified holomorphic
structure.

For the other type of brane, with $T'\lambda|=-\lambda|$,
supersymmetry requires $\{\G^{IJ}F_{IJ},T'\}\epsilon=0$.  This
leads to two conditions.  The first is that $F_{IJ}=0$ if $I,J$
are in the subspace perpendicular to directions $1,5,2-7,$ and
$3+6$.  In particular, $F_{0I}=0$ if direction $I$ is in this
subspace, and we impose Dirichlet boundary conditions for such
$I$. Thus, in particular, $A_2+\phi_3$ and $A_3-\phi_2$ obey
Dirichlet boundary conditions. The second consequence of
$\{\G^{IJ}F_{IJ},T'\}\epsilon=0$ is that $F$ must be anti-selfdual
when restricted to the subspace generated by directions
$1,5,2-7,3+6$. This leads to modified Neumann boundary conditions
on $A_5,$ $A_{2-7}$, and $A_{3+6}$: % 式 (12.29)
\setcounter{equation}{28}
\begin{equation}\label{eq:nono}
\begin{aligned}
F_{15}+F_{2-7,3+6}&=0\\ F_{1,2-7}+F_{3+6,5} & =0\\
F_{1,3+6}+F_{5,2-7}&=0.\\\end{aligned}
\end{equation}

For this type of brane, $A_{2+7}$ and $A_{3-6}$ are fixed on the
boundary while $A_{2-7}$ and $A_{3+6}$ are free to fluctuate.  So,
from a low energy point of view, the roles of $\CA_z$ and
$\CA_{\bar z}$ are exchanged, relative to the branes with
$T'\lambda|=\lambda|$.  Thus we get a brane supported on a
submanifold $W''$ of $\MH$ that parametrizes flat $G_{\Bbb{C}}$
bundles with a fixed antiholomorphic structure.  This again is a
brane of type $(A,B,A)$.

There is another and more trivial way to exchange the roles of
$\CA_z$ and $\CA_{\bar z}$: we could reverse the sign of $\phi_2$
and $\phi_3$ or in other words replace $T'$ by
$T''=-\Gamma_{15}\Gamma_{2+7}\Gamma_{3-6}$.  So we really have two
families of branes supported on manifolds parameterizing bundles
with fixed holomorphic or antiholomorphic structure.  The two
families differ by the conditions placed on modes that are massive
away from singularities of $\MH$; they may differ at singularities
of $\MH$.

\par\medskip{\itshape $S$-Duality And The Canonical Coisotropic Brane}

By applying $S$-duality to the examples that we just constructed,
we can obtain new branes, which will also be branes of type
$(A,B,A)$. We know how to apply $S$-duality to the condition for
unbroken supersymmetry: $T'\epsilon=\epsilon$ is replaced by $\hat
T\epsilon=\epsilon$ where % 式 (12.30)
\setcounter{equation}{29}
\begin{equation}\label{eq:nofo}
\hat T ={1+\HG\over \sqrt
2}T'{1-\HG\over \sqrt 2}=-{\Gamma_{15}\over
2}\left((\Gamma_{26}+\Gamma_{37})-\HG(\Gamma_{23}+\Gamma_{67})\right).
\end{equation}

To get new $(A,B,A)$-branes in this way, we must construct
boundary conditions that preserve the supersymmetries with $\hat
T\epsilon=\Gamma_{2367}\epsilon=\epsilon$.

We look for a supersymmetric brane with $\hat
T\lambda|=-\lambda|$. After a somewhat lengthy calculation, we
deduce from the condition $\{\G^{IJ}F_{IJ},\hat T\}\epsilon=0$ for
supersymmetry that the boundary conditions on the bosons must be
% 式 (12.31)
\setcounter{equation}{30}
\begin{equation}\label{eq:tomgo}
\begin{aligned}F_{15}+F_{26}+F_{37}& = 0 \\
                    F_{12}+F_{06}-F_{56}& = 0 \\
                    F_{13}+F_{07}-F_{57}& = 0 \\
                    F_{16}+F_{52}-F_{02}& = 0 \\
                    F_{17}+F_{53}-F_{03}& = 0 ,\\\end{aligned}
\end{equation}

along with $F_{0J}=0$ for $J=4,8,9$.

We set the massive fields $A_J,\,J=4,8,9$ to zero on the boundary.
Otherwise, all fields obey modified Neumann boundary conditions
and are free to fluctuate on the boundary, so the $(A,B,A)$-brane
obtained this way is a space-filling brane whose target space is
all of $\MH$.

We have seen mixed Neumann-Dirichlet boundary conditions before,
but there is a key difference here. In previous examples, the
boundary conditions become purely Dirichlet or purely Neumann if
one sets to zero the fields $A_4,A_5,A_8$, and $A_9$ that play no
role in the low energy theory away from singularities, so they are
irrelevant in the effective two-dimensional field theory of the
brane.  In the present example, this is not the case. After
discarding the ``massive'' fields and choosing the gauge $A_1=0$,
the boundary conditions become % 式 (12.32)
\setcounter{equation}{31}
\begin{equation}\label{eq:bdrycc}
\begin{aligned}
D_0\phi_2+\partial_1 A_2&=0,\\ D_0\phi_3+\partial_1 A_3&=0,\\
F_{02}-\partial_1\phi_2&=0,\\ F_{03}-\partial_1\phi_3&=0.\end{aligned}
\end{equation}

Mixed Dirichlet-Neumann boundary conditions on the sigma-model
fields arise when the gauge field on the brane is not flat;
specifically, the boundary condition is % 式 (12.33)
\setcounter{equation}{32}
\begin{equation}\label{eq:zembro}
f_{MN}\partial_0 X^N+g_{MN}\partial_1 X^N=0,
\end{equation}
 where $f$ is the
curvature of the target-space gauge field, and $g$ is the
target-space metric. Comparing with \eqref{eq:bdrycc}\ we find % 式 (12.34)
\setcounter{equation}{33}
\begin{equation}\label{eq:membro}
f={1\over 2\pi}\int_C |d^2z| \Tr\, \left(\delta \phi_z\wedge\delta
A_\bz+\delta\phi_\bz\wedge \delta A_z\right)=\omega_J.
\end{equation}
 This is
the target-space gauge field for the canonical coisotropic
$(A,B,A)$-brane in the sigma model with target $\MH(G,C)$. So that
is the brane that we have found.

\par\medskip{\it Boundary Observables for the Canonical Coisotropic Brane}

Now we want to make a simple observation about boundary
observables for the canonical coisotropic brane.  We consider
local operators ${\EUO}$ inserted at a point $P$ of the boundary
of $M$. The boundary, of course, looks like $\Bbb{R}\times C$.  We
will consider the cohomology of the topological supercharge $Q$ of
the $A$-model in complex structure $K$ acting on local operators
inserted at a point $s\times p\in \Bbb{R}\times C$.  The metric of
$C$ is irrelevant mod $\{Q,\dots\}$, so we can assume $C$ to be
flat near $p$.

We can classify local operators by their dimension and also by the
``spin'' with which they transform under rotations of $C$ around
the point $p$.  (The considerations will be local, so it will not
matter that the global structure of $C$ typically breaks the
rotation symmetry.)  The natural dimensions in the twisted theory
are 1 for $A,\phi,\psi,\tilde\psi$, and 2 for
$\chi,\eta,\tilde\eta$.

We consider $Q$-invariant operators of dimension $n$ and spin $n$.
To have these quantum numbers, a gauge-invariant operator must be
constructed only from $\phi_z$, $\psi_z$, $\tilde\psi_z$, and the
covariant derivative $D_z$.  But for an operator to be
$Q$-invariant, $D_z$ cannot appear, since $A_z$ is not
$Q$-invariant even on the boundary.  The boundary condition $\hat
T\psi|=-\psi|$ implies that
% 式 (12.35)
\setcounter{equation}{34}
\begin{equation}\label{eq:bdrry}
\begin{aligned}\psi_2|&=-i\tilde\psi_3|\\
\psi_3|&=i\tilde\psi_2|.\\\end{aligned}
\end{equation}
 With $\psi_z=(\psi_2-i\psi_3)/2$,
this leads to $\psi_z|=\tilde\psi_z|$. So $Q$-invariant boundary
observables of dimension $n$ and spin $n$ are functions just of
$\phi_z$ and $\psi_z$, without any derivatives.

 The fields $\phi_z$ and $\psi_z$ have ghost number $\EUK=0$ and
 $\EUK=1$, respectively.  So an operator of this type that also has
 $\EUK=0$ is a gauge-invariant function of $\phi_z$ only, such as
 ${\EUO}=\Tr\,\phi_z^n$ for some $n$.  A typical operator of
 $\EUK=1$ is ${\EUO}'={\rm Tr}\,\phi_z^{n-1}\psi_z$.

 Classically, the boundary conditions ensure that on the boundary
 $[Q,\phi_z\,]=\{Q,\psi_z\,\}=0$, so all operators ${\EUO}$
 and ${\EUO}'$ of the sort just described are nonzero elements
 of the cohomology of $Q$.  What happens quantum mechanically?
 Could there be a quantum correction to the action of $Q$, such that the
 exact quantum formula would be $[Q,{\EUO}]=\epsilon{\EUO}'$
 for some $\epsilon$?  If so, at the quantum level, ${\EUO}$ and
 ${\EUO}'$ would pair up and disappear from the cohomology.

 That this does not occur can be argued using time-reversal
 symmetry.  We define an orientation-reversing symmetry ${\EUT}$
 that reverses the sign of the time coordinate $x^0$ and acts
 trivially on other coordinates of $M$.  Allowing for the minus
 sign with which ${\EUT}$ acts on $\phi$ (see eqn. \eqref{eq:yoro}), ${\EUT}$
 is a symmetry of the boundary conditions that define the
 canonical coisotropic brane.  ${\EUT}$ commutes with the
 topological supercharge of the $A$-model in complex structure
 $K$, since $t=1$ is a fixed point of ${\EUT}$ (see eqn. \eqref{eq:zongor}).
 Also, $\phi_z$ is odd under ${\EUT}$, in view of \eqref{eq:yoro}, but
 \eqref{eq:zoro}\
 plus the boundary condition $\psi_z|=\tilde\psi_z|$ implies that
 $\psi_z|$ is even under ${\EUT}$.  So the operators ${\EUO} $
 and ${\EUO}'$ cannot pair up under the action of $Q$, and must
 survive in the cohomology.   More generally, consider a local
 boundary operator of dimension and spin $n$ and ghost number $k$
 that is $Q$-invariant at the classical level.  Such an operator
 is determined by a gauge-invariant function $F(\phi_z,\psi_z)$
 that is of degree $n-k$ in $\phi_z$ and degree $k$ in $\psi_z$.
 It transforms as $(-1)^{n-k}$ under ${\EUT}$.   For $Q$ to act
 non-trivially on such an operator at the quantum level, it would
 have to leave $n$ unchanged and increase $k$ by 1.  This would
reverse the eigenvalue of ${\EUT}$, and so is impossible.  Thus,
the cohomology of $Q$ in the space of local operators of dimension
and spin $n$ coincides with the classical result, and is given
simply by the space of gauge-invariant functions
$F(\phi_z,\psi_z)$.
